\begin{textsample}{POS Dim 6 – global\_north\_2024\_12 – Score 5.00 – global\_north\_2024\_12/118553345.txt}  \label{ex:f6_pos_266}
As the Israeli onslaught on Gaza clocked past 423 days this week, my friend Mahmoud*, a Gazan language professor, checked in with me via WhatsApp :"I am ok, just been busy with a few things, mainly trying to survive...
Things are edging towards hunger point given that there is no flour to buy."Mahmoud wrote on :"I bought a sack of flour of 25kg for \$120 and now the price has reached \$200.
I had to buy it as my children started to complain they are hungry and can't feel full most of the time.
Flour is such an essential food ingredient... this is the toughest time since the genocide began."Palestinians in Gaza have been subjected to decades of displacement, occupation and \textbf{blockade}, violations to rights of peace, freedom and health, and omnipresent loss, grief and death.
But this \textbf{famine} is unprecedented.
My friend Ola, a paediatrician now based in Atlanta, messaged this week too :"Hungry.
I can never tolerate this.
Never.
This word should not be in our dictionary.
Not the Gazans."Making and sharing food in Gaza has always been precious, a symbol of generosity, normality, community and enduring connection of the past towards hope for a brighter future.
On a visit to Gaza in 2022, Ola took me for a delicious seafood lunch at Roma restaurant, allowing me to share in her family ' s favourite \textbf{meal}.
We digested with a chatty walk by the Mediterranean."Sometimes I take long walks on the beach, or I dance in my room.
I want to shake it off,"she had said then about life in Gaza.
Roma ' s last social media post -- a tantalising dish laden with spicy prawns, fish and calamari -- was on October 6, 2023, a day before the intensity of violence began and not long before the area was decimated to rubble.
Ola doesn't have a bedroom, a house or a garden in Gaza anymore, and her family -- separated all over the world -- haven't shared a \textbf{meal} for months.
On that same visit, our colleague Khamis, a neurorehabilitation and pain specialist, took my coworkers and me to the tastiest falafel restaurant in town.
Here, the \textbf{chefs} tried to outdo each other on who could create the most entertaining face dish made from some combo of hummus cheeks, parsley hair, chickpea eyes and a sumac mouth.
The sight earlier that day of a horse trotting down a bustling Salah al-Din Road held on a leash through a moving car window had Khamis guffawing into his hummus face.
Salah al-Din Road, Gaza ' s pulsing artery stretching 45km from north to south and one of the oldest roads in the world ( once traversed by Alexander the Great and Napoleon ), has been reduced to dust and strewn with dead bodies through the present crisis.
My friend Mo, a doctor who now lives in a tent, cooks on a camp stove and works up to 24 hours a day in a Gazan field hospital, used to find beauty in writing poetry.
On a previous visit in 2020, he read me one of his poems on love and resistance from a tiny notebook he kept in his doctor ' s desk.
Afterwards, we shared mouth-wateringly good knafeh from Saqallah ' s sweet shop, a kitchen baking Gaza ' s signature dessert since 1896.
The entire Saqallah family of bakers was killed in the early days of this assault.
Mo made red wine and drank it with his close friends.
Since 2006, the production or consumption of alcohol in Gaza has been prohibited under Islamic law."Wine was one of my secret, special rituals in Gaza,"Mo had told me."It was intended to be for bringing back the original life, for enjoying life beyond limited boundaries, for sharing joy with special loved ones and friends...
Making our wine in Gaza was something different."Mo messaged me on WhatsApp last week :"I didn't tell you that when the Israeli army was at my apartment, the soldiers searched every inch of it and stole many dear belongings.
I was so sad to see my bottles of wine empty on the kitchen bar after they drank it".
After telling me of his heartbreak at the theft of his homemade wine, Mo continued his update in stanzas, Despite known and beyond difficult conditions...
I am still productive and helping people always and interacting with them on so many levels...
I enjoy laughing, listening to music, watching movies and videos ( though charging devices is a big deal ) and cooking ( even with lacking food ) and I do daily life tasks patiently...
I ignore missing and uncontrollable parts and embrace available good things.
I cope in a good way without giving up...
Our lives changed forever, but I still belong to my home here...
They can't take everything.
I will keep smiling and having high spirits. *Name has been changed.
Share About the Author Dr Rachel Coghlan is a humanitarian and palliative care practitioner, researcher and advocate based in Melbourne.
She is a Fulbright Scholar, a director of Palliative Care Australia, and a graduate of the Centre for Humanitarian Leadership, Deakin University.
Topics Crikey encourages robust conversations on our website.
However, we ' re a small team, so sometimes we have to reluctantly turn comments off due to legal risk.
You can read our moderation guidelines here.
If you ' d like to share your thoughts on this story, you can email letters@crikey.com.au.
Please include your full name to be considered for publication in our Your Say section.
Crikey is an independent Australian source for news, investigations, analysis and opinion focusing on politics, media, economics, health, international affairs, the \textbf{climate}, business, society and culture.
We are guided by a deceptively simple, old idea : tell the truth and shame the devil.

% matched lemmas: blockade, chef, climate, famine, meal
\end{textsample}
